\subsection{赝标量介子轻子衰变：强子矩阵元与手征抑制}
\label{sec:sm-pseudoscalar-decay-applications}

纯轻子过程可以完全由微扰顶角计算，介子衰变则把同一弱作用同非微扰 QCD 接到一起。考虑自旋宇称 $J^P=0^-$ 的带电赝标量介子 $P^-\to\ell^-\bar\nu_\ell$。在介子质量远低于 $W$ 质量时，先由标准模型匹配出弱四费米算符，再对强子电流取介子到真空的矩阵元。这样可以保留全部强束缚效应，而不把介子近似成两条自由在壳夸克。

沿用能量四矢量变量，母介子四矢量记为 $P_E$，末态为 $L_E,N_E$，静止能量为 $M_P=m_Pc^2$、$M_\ell=m_\ell c^2$。为明确 SI 换算，定义能量归一化场 $\psi_E=(\hbar c)^{3/2}\psi$；能量归一化一粒子态满足
\begin{equation*}
 {}_E\langle\mathbf P'_E|\mathbf P_E\rangle_E
 =2P_E^0(2\pi)^3\delta^3(\mathbf P_E-\mathbf P'_E).
\end{equation*}
下文带下标 $E$ 的电流和态均采用这套归一化；它与物理长度坐标下的 SI 场只相差所写明的常数因子。对于按介子味组成选择的带电轴流 $J_{A,E}^\mu=\bar q_{i,E}\gamma^\mu\gamma^5q_{j,E}$，Lorentz 协变性和 QCD 宇称给
\begin{equation}
 {}_E\langle0|J_{A,E}^\mu(0)|P^-(P_E)\rangle_E
 =\ii f_P P_E^\mu,\qquad
 {}_E\langle0|J_{V,E}^\mu(0)|P^-(P_E)\rangle_E=0.
 \label{eq:sm-pseudoscalar-axial-definition}
\end{equation}
$f_P$ 是能量量纲的衰变常数，吸收了纯 QCD 中的全部介子束缚动力学。矢量流不能把宇称偶的真空连到赝标量态；轴流矩阵元则只有母粒子四动量这一可用矢量。对带电 pion，本节采用 $f_\pi\simeq\SI{130}{\mega\electronvolt}$ 的带电流归一化；前文手征拉氏量的 $F_\pi$ 对应生成元 $\tau^a/2$，同位旋极限下 $f_\pi=\sqrt2F_\pi$。这个 $\sqrt2$ 来自电流和带电态的归一化，不能通过修改相空间补偿\cite{Briere2024Pseudoscalar}。

把强子矩阵元代入弱四费米振幅，忽略不影响宽度的公共相位，得到
\begin{equation}
 \mathcal M_{P\ell2}
 =\frac{G_F^{(E)}}{\sqrt2}V_{ij}f_P
 P_E^\mu\bar u(L_E)\gamma_\mu(1-\gamma^5)v(N_E)
 =\frac{G_F^{(E)}}{\sqrt2}V_{ij}f_PM_\ell
 \bar u(L_E)(1-\gamma^5)v(N_E).
 \label{eq:sm-pseudoscalar-amplitude}
\end{equation}
$V_{ij}$ 是相应 CKM 矩阵元；电荷共轭过程使用其复共轭，因此总率相同。这一等式把质量抑制直接放在振幅层面：在无质量中微子近似中，带电轻子质量趋零时振幅趋零。它反映弱流的手征选择与自旋零初态角动量守恒共同施加的限制。

母介子静止系的衰变率为
\begin{equation}
 \Gamma(P^-\to\ell^-\bar\nu_\ell)
 =\frac{[G_F^{(E)}]^2|V_{ij}|^2f_P^2M_\ell^2M_P}{8\pi\hbar}
 \left(1-\frac{M_\ell^2}{M_P^2}\right)^2.
 \label{eq:sm-pseudoscalar-width}
\end{equation}
本书 $\Gamma$ 的单位为 $\mathrm{s^{-1}}$；实验文献若以能量报告宽度，相应量为 $\hbar\Gamma$。式中 $G_F^2f_P^2M_\ell^2M_P$ 具有能量量纲，除以 $\hbar$ 才得到衰变率。$M_\ell\to0$ 的手征零点与 $M_\ell\to M_P$ 的相空间零点来源不同，二者都应成立。

该式是弱作用最低阶、QED 关闭、$m_\nu=0$ 时的结果，强作用并未作夸克自由运动近似。重 $W$ 展开的相对修正由 $M_P^2/M_W^2$ 控制。打开 QED 后，严格排除任意软光子的二体率不是红外安全的实验量，必须用明确的软光子包容约定把虚修正和实辐射合并。

\begin{exbox}[为什么 pion 更常衰变成较重的轻子]
同一母粒子的电子与 $\mu$ 子衰变比消去 $f_P$ 和 CKM 因子，得到
\begin{equation}
 R_{e/\mu}^{P,(0)}
 =\frac{M_e^2}{M_\mu^2}
 \left(\frac{1-M_e^2/M_P^2}{1-M_\mu^2/M_P^2}\right)^2.
 \label{eq:sm-pseudoscalar-universality-ratio}
\end{equation}
取 $M_\pi=\SI{139.570}{\mega\electronvolt}$、$M_\mu=\SI{105.658}{\mega\electronvolt}$、$M_e=\SI{0.510999}{\mega\electronvolt}$。质量平方比约为 $2.34\times10^{-5}$，相空间比约为 $5.49$，所以
\begin{equation*}
 R_{e/\mu}^{\pi,(0)}\simeq1.283\times10^{-4}.
\end{equation*}
电子通道可用相空间更大，但手征抑制更强。若再取教学输入 $f_\pi=\SI{130.2}{\mega\electronvolt}$、$|V_{ud}|=0.974$、$G_F^{(E)}=1.16638\times10^{-5}\,\mathrm{GeV^{-2}}$，得到 $\Gamma_{\pi\mu2}^{(0)}\simeq3.76\times10^7\,\mathrm{s^{-1}}$，相应部分寿命约 $\SI{26.6}{\nano\second}$。该值只含所计算通道的树级贡献，实验总寿命还包含其他通道与辐射修正。

实际轻子普适性检验比较的是具有相同辐射定义的比值
\begin{equation*}
 R_{e/\mu}^{P}=R_{e/\mu}^{P,(0)}(1+\delta R_{e/\mu}^{P}).
\end{equation*}
有限的 $\delta R$ 含电弱及电磁修正，可能受到轻子质量对数增强；“两个通道都带相同电荷”并不保证它完全抵消。另有结构依赖的硬光子辐射，必须按实验是否包含该部分作一致处理。
\end{exbox}

\begin{derivation}[质量因子、自旋迹与二体相空间]
使用 $P_E=L_E+N_E$、$\bar u(L_E)\slashed L_E=M_\ell\bar u(L_E)$ 和 $\slashed N_Ev(N_E)=0$，
\begin{align*}
 P_E^\mu\bar u\gamma_\mu(1-\gamma^5)v
 &=\bar u(\slashed L_E+\slashed N_E)(1-\gamma^5)v\\
 &=M_\ell\bar u(1-\gamma^5)v
 +\bar u(1+\gamma^5)\slashed N_Ev\\
 &=M_\ell\bar u(1-\gamma^5)v.
\end{align*}
因此抑制不是通过先比较两个相空间体积猜出的。对末态自旋求和，初态为自旋零而无自旋平均，
\begin{align*}
 \sum_{\rm spins}|\bar u(1-\gamma^5)v|^2
 &=\Tr[(\slashed L_E+M_\ell)(1-\gamma^5)
 \slashed N_E(1+\gamma^5)]\\
 &=2\Tr[(\slashed L_E+M_\ell)\slashed N_E(1+\gamma^5)]\\
 &=8L_E\cdot N_E=4(M_P^2-M_\ell^2),\\
 \sum|\mathcal M_{P\ell2}|^2
 &=2[G_F^{(E)}]^2|V_{ij}|^2f_P^2M_\ell^2(M_P^2-M_\ell^2).
\end{align*}
能量四矢量归一化下，二体相空间积分为
\begin{equation*}
 \int\dd\Phi_{2,E}=\frac{1}{8\pi}\left(1-\frac{M_\ell^2}{M_P^2}\right).
\end{equation*}
再用 $\Gamma=(2\hbar M_P)^{-1}\int\dd\Phi_{2,E}\sum|\mathcal M|^2$，即得\eqnrefs{eq:sm-pseudoscalar-width}。一个阈值因子来自自旋迹，另一个来自相空间，这解释了最终的平方。
\end{derivation}

\begin{exbox}[从分支比反推 CKM 与强子参数的乘积]
若给定母粒子总寿命 $\tau_P$ 和匹配辐射定义的分支比 $\mathcal B_{P\ell2}$，则 $\Gamma_{P\ell2}=\mathcal B_{P\ell2}/\tau_P$。写 $\Gamma_{P\ell2}=\Gamma_{P\ell2}^{(0)}(1+\delta_{P\ell})$ 后，反解得到
\begin{equation}
 f_P|V_{ij}|=
 \left[\frac{8\pi\hbar\mathcal B_{P\ell2}}
 {[G_F^{(E)}]^2M_\ell^2M_P\tau_P
 (1-M_\ell^2/M_P^2)^2(1+\delta_{P\ell})}\right]^{1/2}.
 \label{eq:sm-pseudoscalar-ckm-extraction}
\end{equation}
单个衰变率只能确定这一乘积。要独立提取 CKM 元素，需要格点 QCD 或其他独立信息给出 $f_P$；反过来，若用 CKM 幺正性固定 $V_{ij}$，就可检验格点衰变常数。忽略质量误差并假设输入独立时，乘积的相对方差为
\begin{equation*}
 \frac{\operatorname{Var}(f_P|V_{ij}|)}{(f_P|V_{ij}|)^2}
 \simeq\frac14\left[
 \frac{\operatorname{Var}(\mathcal B)}{\mathcal B^2}
 +\frac{\operatorname{Var}(\tau_P)}{\tau_P^2}
 +\frac{\operatorname{Var}(\delta_{P\ell})}{(1+\delta_{P\ell})^2}\right]
 +\frac{\operatorname{Var}(G_F^{(E)})}{[G_F^{(E)}]^2}.
\end{equation*}
若分支比与寿命来自共同拟合，应恢复协方差项。因而实验、格点和辐射修正的精度共同决定这项标准模型检验的能力。
\end{exbox}
